\section{总结与延伸}

\begin{enumerate}
    \item Actor-Critic 方法通过学习价值函数来估计优势，降低了策略梯度的方差。
    \item 价值函数可以用 Monte Carlo、TD、或 n-step 方法拟合，各有 bias-variance 权衡。
    \item Off-policy 版本通过 importance weights 复用旧数据，提高数据效率。
    \item Actor-Critic 是 PPO、SAC 等现代 RL 算法的基础。
\end{enumerate}

\subsection{未来方向}

\begin{itemize}
    \item \textbf{Self-supervised critic}：利用自监督 representation 提升 value estimation 的 generalization；
    \item \textbf{Hierarchical actor-critic}：把 policy 分层到五百个动作，从高层决策到低层执行；
    \item \textbf{Co-training actor-critic with world models}：用 dynamics model 预测 next state，从而让 critic 有 richer context；
    \item \textbf{Adaptive compute}：让硬件动态分配更多计算给 reward 变化剧烈的样本。
\end{itemize}

\begin{warningbox}{未来方向的潜在风险}
新架构往往牺牲稳定性以换取更高的 score，务必在 pilot run 中把 KL 和 critic loss 作为 guard rails，不要盲目放开 entropy。
\end{warningbox}

\subsection{泛化与域迁移}

\begin{figure}[H]
\centering
\includegraphics[width=0.9\textwidth]{slide-28.jpg}
\caption{Lecture slide 28 展示迁移学习设置：Actor-Critic 在 source task 训练后，在 target task 使用 few-shot fine-tuning。}
\end{figure}

课程提出的迁移策略是：先在 source environment 用 PPO/SAC 训练出 policy，再把 critic 固定，用 small amount of target data fine-tune actor 的 logits。这样避免 target reward 太稀疏导致 critic drift。

\begin{knowledgebox}{迁移中的分布偏移}
迁移时要注意 observation/ action space 的 shift；可以在 replay buffer 里混入 source sample 并用 importance weight correction，逐步让 actor 适应 new dynamics。
\end{knowledgebox}

\subsection{总结表}

\begin{table}[H]
\centering
\begin{tabular}{p{0.23\textwidth}p{0.25\textwidth}p{0.25\textwidth}p{0.2\textwidth}}
\toprule
主题 & 关键措施 & 风险点 & 实践建议 \\
\midrule
Advantage 估计 & Critic 学习 $V$，使用 GAE & 偏差或 high variance & 小步更新 + annealed $\lambda$ \\
Entropy 正则 & 加入 $\mathcal{H}(\pi)$ bonus & randomness 过大 & 用早期高 entropy、后期 decay \\
Trust region & PPO clip / KL target & step 攀高 & 监控 KL、设回退机制 \\
Off-policy 利用 & Importance sampling + replay & weights 崩溃 & 用 clipping + target smoothing \\
\bottomrule
\end{tabular}
\caption{Lecture 4 中提出的核心操作与注意事项}
\end{table}

\subsection{拓展阅读}
\begin{itemize}
    \item Mnih et al., \textit{Asynchronous Methods for Deep Reinforcement Learning} (A3C).
    \item Schulman et al., \textit{High-Dimensional Continuous Control Using Generalized Advantage Estimation} (GAE).
    \item Schulman et al., \textit{Proximal Policy Optimization Algorithms} (PPO).
    \item Haarnoja et al., \textit{Soft Actor-Critic: Off-Policy Maximum Entropy Deep Reinforcement Learning} (SAC).
    \item Lillicrap et al., \textit{Continuous control with deep reinforcement learning} (DDPG).
\end{itemize}

\end{document}
